% 第4節 選択肢集合の困難とネットワーク型選択モデル
%% chapters/behavior_modeling/ch3_choiceset.tex
%% 第3章　選択肢集合
\section{選択肢集合の困難とネットワーク型選択モデル}\label{sec:choiceset}
本節では，選択肢集合の問題と経路非列挙型経路選択モデルについて解説する．
具体的には，選択肢集合の定義の難しさと列挙型モデルの限界を説明した上で，Recursive LogitやDynamic Decision Modelなどの経路非列挙型の経路選択モデルを紹介する．

\subsection{選択肢集合の爆発}


% 4.1 選択肢集合列挙問題
% 経路は膨大
% 活動系列はさらに膨大
% choice set の定義自体が難しい
% 4.2 列挙型経路選択モデルの限界
% サンプリングバイアス
% 列挙ルール依存
% 観測されない代替経路
% 4.3 Dynamic / Recursiveな表現
% Recursive Logit
% Dynamic decision model
% Perturbed Utility Route Choice
% リンク単位・状態単位での逐次選択

% 4.4 
% 多主体相互作用・時空間制約

% この節はかなり重要です。
% 「交通分野では経路選択を普通のMNLに入れれば済むわけではない」という話が出せるからです。